The stack is the one a quantitative desk uses:
\codehl{numpy} and \codehl{scipy} for the numerics,
\codehl{pandas} and \codehl{pyarrow} for tabular data,
\codehl{pandera} for the schemas that say what a table is,
\codehl{scikit-learn} for the fitting,
and \codehl{matplotlib} and \codehl{plotly} for the figures.

Four habits are taught alongside them.

A type owns its own serialisation:
the schema that validates it,
the pair of functions that write it to a table and read it back,
and a test that the round trip is the identity.

A parametrisation is a frozen object and never a computed default.
A default invites a reader to skip the choice,
and in a model the choice is usually the subject.

The clear implementation comes first and is understood before anything faster is written.
An optimisation is then a second implementation shown to beat the first,
and the comparison is part of the material.

A claim about what code does is written as a test,
so that it is re-run whenever the code changes.

Work is kept under version control,
and reaches the course repository through branches and pull requests,
with the tests run automatically on each.
Section \ref{sec.assignments} uses that route,
and it is part of what the assignments teach.
